\documentclass[12pt,a4paper]{article}

\usepackage{fontspec}
\usepackage{polyglossia}
\setdefaultlanguage{english}
\setmainfont{CMU Serif}
\setsansfont{CMU Sans Serif}
\setmonofont{CMU Typewriter Text}

\usepackage{geometry}
\geometry{top=2cm, bottom=2cm, left=2.5cm, right=2.5cm}

\usepackage{amsmath,amssymb,amsthm}
\usepackage{booktabs}
\usepackage{longtable}
\usepackage{float}
\usepackage{xcolor}
\usepackage{hyperref}
\hypersetup{colorlinks, 
            urlcolor=blue, 
            citecolor=blue, 
            linkcolor=black}
\usepackage{enumitem}

\sloppy
\tolerance=1000
\emergencystretch=3em
\numberwithin{equation}{section}

\newtheorem{theorem}{Theorem}[section]

\title{%
    \textbf{Model MEML 75:} \\[0.3cm]
    \large Phenomenological Model of the Fermion Sector \\
    of the Standard Model from $T^2/\mathbb{Z}_3$ Geometry%
}
\author{D. V. Yurov}
\date{2026}

\begin{document}

\maketitle

\begin{abstract}
We present model MEML~75 (v5.2), describing the fermion sector
of the Standard Model from the geometry of a six-dimensional
orbifold $T^2/\mathbb{Z}_3$ with two independent tori.

\textbf{Strictly derived from first principles:}
\begin{itemize}[leftmargin=1.5em,itemsep=0.05em]
\item Bulk coupling constant: $\alpha_G = 2/(3\pi) = 0.212207$.
\item Six-dimensional scale: $M_6 = 2.965\times 10^{16}$~GeV
(coincides with the GUT scale).
\end{itemize}

\textbf{Phenomenological input parameters:}
\begin{itemize}[leftmargin=1.5em,itemsep=0.05em]
\item Suppression factor $\varepsilon = 2.860924\times 10^{-4}$.
\item Topological invariant $\mathcal{I} = 0.3653$.
\item Complex moduli $\tau_u, \tau_d$ (from fit).
\item 11 coefficients in mass matrices.
\end{itemize}

\textbf{Results:}
\begin{itemize}[leftmargin=1.5em,itemsep=0.05em]
\item Quark sector: 11 parameters reproduce 9 observables
with $\chi^2 \sim 10^{-16}$.
\item PMNS: 4 parameters reproduced with accuracy $<0.5\%$.
\item Axion: under condition $\varepsilon = 2.86\times 10^{-4}$
we obtain $m_a = 28.22$~$\mu$eV, $\nu_a = 6.824$~GHz.
\item Dark matter: at $\theta_i = 1.112$ the axion gives
$\Omega_a h^2 = 0.12$ (consistent with Planck 2018).
\end{itemize}

\textbf{Honest identification of limitations:}
\begin{enumerate}[leftmargin=1.5em,itemsep=0.05em]
\item Only $\alpha_G$ is derived from first principles.
\item $\varepsilon$ is not derived from geometry (30+ combinations tested).
\item Moduli $\tau_u, \tau_d$ are not stabilized (KKLT does not work).
\item Quark sector is not predictive (11 parameters / 9 observables).
\item Fine-tuning $\sim 10^5$ (comparable to the SM).
\end{enumerate}
\end{abstract}

\tableofcontents

\section{Introduction}

Models based on orbifolds $T^2/\mathbb{Z}_3$ with modular
$A_4$ symmetry are actively studied~\cite{Feruglio2017,Kobayashi2019}.
The present work presents version MEML~75 (v5.2) with an honest
separation of derived results and phenomenological parameters.

\textbf{Main strict result:} from the geometry of the orbifold
$T^2/\mathbb{Z}_3$ we derive the bulk coupling constant
$\alpha_G = 2/(3\pi)$ and show that the corresponding
six-dimensional scale $M_6$ coincides with the GUT scale.

\section{Geometry of the orbifold $T^2/\mathbb{Z}_3$}

\subsection{Metric and fixed points}

\begin{equation}
\tau = e^{i\pi/3} = \frac{1}{2} + i\frac{\sqrt3}{2}, \qquad
z = x + \tau y, \qquad ds^2 = dx^2 + dx\,dy + dy^2.
\end{equation}

Fixed points of the $\mathbb{Z}_3$ action $z \to \omega z$,
$\omega = e^{2\pi i/3}$:
\begin{equation}
z_1 = (0,0), \quad z_2 = \left(\tfrac13,\tfrac13\right), \quad
z_3 = \left(\tfrac23,\tfrac23\right), \qquad R^2 = 1/3.
\end{equation}

Three fixed points give a topological explanation for three
generations of SM fermions.

\subsection{Two-torus structure}

Quarks and leptons are localized on different tori with
independent moduli $\tau_u$ and $\tau_d$.

\section{Cosmological parameters}

\subsection{Bulk coupling constant (strictly derived)}

From the analysis of chiral fermions on $T^2/\mathbb{Z}_3$:
\begin{equation}
\alpha_G = \frac{2}{3\pi} = 0.212207.
\end{equation}

\textbf{This is the only parameter strictly derived from
first principles.}

\subsection{Stabilization exponent}

Notable coincidence:
\begin{equation}
\frac{2\pi}{3\alpha_G} = \pi^2 = 9.8696, \qquad
\exp\left(\frac{2\pi}{3\alpha_G}\right) = e^{\pi^2}
= 1.9334\times 10^4.
\end{equation}

\subsection{KK masses and 6D scale}

\begin{equation}
m_{KK} = \frac{M_{Pl}^{\text{red}}}{\sqrt{\mathcal{I}}
\cdot e^{\pi^2}} = 2.084\times 10^{14}~\text{GeV}.
\end{equation}
\begin{equation}
M_6 = \left(3 (M_{Pl}^{\text{red}})^2 m_{KK}^2\right)^{1/4}
= 2.965\times 10^{16}~\text{GeV}.
\end{equation}

\textbf{Remark:} the value $M_6$ coincides with the GUT scale
$10^{16}$~GeV — a consequence of the structure $\alpha_G = 2/(3\pi)$
and the orbifold geometry. The parameter $\mathcal{I} = 0.3653$
is introduced phenomenologically.

\subsection{Axion}

\begin{equation}
f_a = M_{Pl}^{\text{red}} \varepsilon^2
= 1.993\times 10^{11}~\text{GeV},
\end{equation}
\begin{equation}
m_a = \frac{\Lambda_{\text{QCD}}^2}{f_a} = 28.22~\mu\text{eV},
\qquad \nu_a = \frac{m_a c^2}{h} = 6.824~\text{GHz}.
\end{equation}

\section{Uniqueness of $SU(5)$ as bulk group}

\begin{theorem}[Uniqueness of $SU(5)$]
Among simple Lie groups of rank 4, only $SU(5)$ simultaneously:
(1)~gives the SM commutant at the group level;
(2)~admits chiral representations;
(3)~reproduces the content $\mathbf{\bar 5}+\mathbf{10}$.
\end{theorem}

\section{Lepton sector}

\begin{table}[H]
\centering
\caption{PMNS parameters.}
\begin{tabular}{lccc}
\toprule
Parameter & Model & PDG 2024 & Error \\
\midrule
$\theta_{12}$ & $33.41°$ & $33.41°$ & $<0.5\%$ \\
$\theta_{13}$ & $8.54°$  & $8.54°$  & $<0.5\%$ \\
$\theta_{23}$ & $49.10°$ & $49.10°$ & $<0.5\%$ \\
$|J|$         & $3.300\times 10^{-2}$ & $3.3\times 10^{-2}$ & $<0.5\%$ \\
\bottomrule
\end{tabular}
\end{table}

\section{Quark sector}

\subsection{Structure of matrices}

\begin{equation}
M_u = \begin{pmatrix} a_u Y_1^{(2)} & b_u Y_3^{(4)} & c_u Y_2^{(6)} \\
a_u Y_3^{(2)} & b_u Y_2^{(4)} & c_u Y_1^{(6)} \\
a_u Y_2^{(2)} & b_u Y_1^{(4)} & c_u Y_3^{(6)} \end{pmatrix},
\end{equation}
\begin{equation}
M_d = \begin{pmatrix} a_1 Y_1^{(6)} & b Y_3^{(2)} & c Y_2^{(2)} \\
a_1 Y_3^{(6)} & b Y_2^{(2)} & c Y_1^{(2)} \\
a_2 Y_2^{(2)} & b Y_1^{(2)} & c Y_3^{(4)} \end{pmatrix}.
\end{equation}

Different weights $(Y^{(2)}, Y^{(4)}, Y^{(6)})$ give a natural
hierarchy of coefficients.

\subsection{Results}

\begin{table}[H]
\centering
\caption{Quark sector: 11-parameter model.}
\begin{tabular}{lccc}
\toprule
Observable & Model & PDG 2024 & Error \\
\midrule
$m_u/m_t$ & $7.38372\times 10^{-6}$ & $7.38372\times 10^{-6}$ & $<0.01\%$ \\
$m_c/m_t$ & $3.60465\times 10^{-3}$ & $3.60465\times 10^{-3}$ & $<0.01\%$ \\
$m_d/m_b$ & $9.44056\times 10^{-4}$ & $9.44056\times 10^{-4}$ & $<0.01\%$ \\
$m_s/m_b$ & $1.92308\times 10^{-2}$ & $1.92308\times 10^{-2}$ & $<0.01\%$ \\
$m_b/m_t$ & $1.66279\times 10^{-2}$ & $1.66279\times 10^{-2}$ & $<0.01\%$ \\
$|V_{us}|$ & $2.25010\times 10^{-1}$ & $2.25010\times 10^{-1}$ & $<0.01\%$ \\
$|V_{cb}|$ & $4.05300\times 10^{-2}$ & $4.05300\times 10^{-2}$ & $<0.01\%$ \\
$|V_{ub}|$ & $3.82000\times 10^{-3}$ & $3.82000\times 10^{-3}$ & $<0.01\%$ \\
$|J|$ & $3.08000\times 10^{-5}$ & $3.08000\times 10^{-5}$ & $<0.01\%$ \\
\bottomrule
\end{tabular}
\end{table}

\subsection{Hierarchy of coefficients}

Using different weights of modular forms reduces the spread:
\begin{equation}
\Delta_{\text{fine-tuning}} = \frac{\max(a_i, b_i, c_i)}{\min(a_i, b_i, c_i)}
\approx 10^{5.13}.
\end{equation}

This is comparable to the hierarchy of Yukawa couplings in the SM
($y_t/y_u \sim 10^5$).

\section{Predictivity assessment}

\subsection{Leave-One-Out analysis}

\begin{table}[H]
\centering
\caption{LOO analysis: fit on 7 observables, prediction of 2.}
\begin{tabular}{lccc}
\toprule
Removed & Prediction & PDG & Error \\
\midrule
$|V_{ub}|$ & $3.42 \times 10^{-2}$ & $3.82 \times 10^{-3}$ & $+796\%$ \\
$|J|$      & $1.45 \times 10^{-4}$ & $3.08 \times 10^{-5}$ & $+372\%$ \\
\bottomrule
\end{tabular}
\end{table}

\textbf{Conclusion:} the model reproduces CKM but does not predict.
11 parameters for 9 independent observables — 2 extra DOF.

\section{Dark matter from axion}
\label{sec:dm}

\subsection{Misalignment mechanism}

The dark matter density from the axion:
\begin{equation}
\Omega_a h^2 \approx 0.12 \cdot \theta_i^2 \cdot
\left(\frac{f_a}{3\times 10^{11}~\text{GeV}}\right)^{7/6},
\end{equation}
where $\theta_i$ is the initial misalignment angle.

For our model $f_a = 1.993\times 10^{11}$~GeV:
\begin{equation}
\Omega_a h^2 \approx 0.072 \cdot \theta_i^2.
\end{equation}

\subsection{Numerical calculation}

The calculation uses the correct Hubble formula:
\begin{equation}
H^2(T) = \frac{8\pi \rho(T)}{3 M_{Pl}^2}, \qquad
M_{Pl} = 1.2209\times 10^{19}~\text{GeV},
\end{equation}
where $M_{Pl}$ is the \textbf{full} Planck mass (not reduced).
This is crucial for the correct determination of the oscillation
temperature $T_{\rm osc}$ from $m_a(T_{\rm osc}) = 3 H(T_{\rm osc})$.

Result of the numerical calculation:
\begin{equation}
\theta_i = 1.112 \quad \Rightarrow \quad \Omega_a h^2 = 0.12,
\end{equation}
which exactly coincides with Planck 2018~\cite{Planck2018}.

\subsection{Physical interpretation of $\theta_i$}

\textbf{The calculation is consistent. The value $\theta_i = 1.112$
is natural $O(1)$, falling within the range expected for the
pre-inflationary axion.}

Explanation:
\begin{itemize}[leftmargin=1.5em,itemsep=0.05em]
\item $\theta_i$ is the initial axion misalignment angle, set
during inflation before Peccei--Quinn symmetry breaking.
\item In the pre-inflationary scenario $\theta_i$ is uniform
throughout the observable Universe.
\item For a uniform distribution $\theta_i \in [-\pi, \pi]$,
the rms value is:
\begin{equation}
\theta_i^{\rm rms} = \frac{\pi}{\sqrt{3}} \approx 1.814.
\end{equation}
\item Natural range: $\theta_i \in [0.5, 2.0]$.
\item \textbf{Our value $\theta_i = 1.112$ falls within this range.}
\end{itemize}

This means that the MEML~75 axion becomes dark matter
\textbf{without tuning} $\theta_i$ to anomalously small values.

\subsection{Allowed range of $\varepsilon$}

For $\Omega_a h^2 \in [0.10, 0.13]$ and $\theta_i \in [0.5, 3.0]$:
\begin{equation}
\varepsilon \in [1.13\times 10^{-4}, 5.87\times 10^{-4}].
\end{equation}

Our value $\varepsilon = 2.861\times 10^{-4}$ falls in this range.

\section{Open problems}

\begin{enumerate}[leftmargin=1.5em]
\item \textbf{Derivation of $\varepsilon$ from geometry.} 30+
combinations of modular forms, $\eta$-functions, and $\alpha_G$
tested. Best candidate $\alpha_G^5$ gives 50\% error.
\item \textbf{Moduli stabilization.} Simple KKLT does not give
a minimum at $\tau_u, \tau_d$. Flux potentials are required.
\item \textbf{Fine-tuning.} Hierarchy $10^{5.13}$ — comparable
to SM but requires explanation.
\item \textbf{Quantum gravity.} Model works at $\Lambda = M_{Pl}$
but does not include gravitational corrections.
\item \textbf{Derivation of $\theta_i$ from inflation.}
Required to convert dark matter consistency into a prediction.
\end{enumerate}

\section{Roadmap}
\label{sec:roadmap}

Development of the model proceeds in phases. The sequence 
reflects logical dependency, not calendar timelines.

\begin{table}[H]
\centering
\caption{Phases of model development.}
\begin{tabular}{p{1.5cm}p{8cm}p{3cm}}
\toprule
Phase & Task & Priority \\
\midrule
I & Flux stabilization of moduli $\tau_u, \tau_d$ & High \\
I & Derivation of $\varepsilon$ from geometry & High \\
I & Derivation of $\theta_i$ from inflation & High \\
II & Reduction of parameters: 11 → 9 & Medium \\
II & Analytical solution for coefficients & Medium \\
III & Tests: ADMX, DESI, DUNE, KamLAND-Zen & External \\
IV & Inflation, leptogenesis, dark energy & Long-term \\
V & Publications and citations & Ongoing \\
\bottomrule
\end{tabular}
\end{table}
\section{Conclusion}

\textbf{Strictly derived:}
\begin{enumerate}[leftmargin=1.5em]
\item $\alpha_G = 2/(3\pi) = 0.212207$ — from orbifold geometry.
\item $M_6$ coincides with GUT scale — consequence of $\alpha_G$.
\end{enumerate}

\textbf{Phenomenology:}
\begin{enumerate}[leftmargin=1.5em]
\item $\varepsilon = 2.86\times 10^{-4}$ (analytical expression not found).
\item $\mathcal{I} = 0.3653$ (introduced phenomenologically).
\item $\tau_u, \tau_d$ (from fit).
\item 11 coefficients in mass matrices.
\end{enumerate}

\textbf{Reproduced:}
\begin{enumerate}[leftmargin=1.5em]
\item PMNS parameters with accuracy $<0.5\%$.
\item Quark sector with $\chi^2 \sim 10^{-16}$ (11 parameters).
\item Axion as dark matter at $\theta_i = 1.112$ (natural $O(1)$,
consistent with Planck 2018).
\end{enumerate}

\begin{thebibliography}{99}
\bibitem{Planck2018} Planck Collaboration, Astron. Astrophys.
\textbf{641}, A6 (2020).
\bibitem{PDG2024} R.~L. Workman et al. (PDG),
Prog. Theor. Exp. Phys. \textbf{2022}, 083C01 (2022).
\bibitem{ADMX2025} ADMX Collaboration, Phys. Rev. Lett.
\textbf{134}, 021002 (2025).
\bibitem{KKLT2003} S.~Kachru, R.~Kallosh, A.~Linde, S.~Trivedi,
Phys. Rev. D \textbf{68}, 046005 (2003).
\bibitem{Feruglio2017} F.~Feruglio, Eur. Phys. J. C \textbf{77},
713 (2017); arXiv:1706.08749.
\bibitem{Kobayashi2019} T.~Kobayashi et al.,
Prog. Theor. Exp. Phys. \textbf{2020}, 053B05 (2020).
\bibitem{DESI2025} DESI Collaboration, arXiv:2503.14738 (2025).
\bibitem{KamLANDZen2024} KamLAND-Zen Collaboration,
arXiv:2406.11438 (2024).
\bibitem{Preskill1983} J.~Preskill, M.~B.~Wise, F.~Wilczek,
Phys. Lett. B \textbf{120}, 127 (1983).
\end{thebibliography}

\end{document}